% =============================================================
% III. THE NEW SERIES
% =============================================================
\section{The new series}
\label{sec:results}

The annual series of GDP per capita and population for 1913 and 1919--2022 is reported in full in Table~\ref{tab:series} of the appendix, and the sectoral volume indices from which it is built in Tables~\ref{tab:appsec1} and~\ref{tab:appsec2}. Table~\ref{tab:growth} summarises growth by sub-period alongside the comparators that recur in the discussion below, and Table~\ref{tab:sectoral} gives the growth of each sector.

% tab:growth
\begin{table}[htbp]
\centering\small\setlength{\tabcolsep}{3.5pt}
\caption[Average annual growth of GDP per capita by sub-period (per cent)]{\textsc{Average annual growth of GDP per capita by sub-period (per cent)}}
\label{tab:growth}
\begin{threeparttable}
\begin{tabular}{lrrrrrrrr}
\toprule
Period & Estonia & Lithuania & Finland & Nordics & Western core & Western periphery & ECE & USSR \\
\midrule
\textit{1913--38} & 1.1 & 0.9 & 2.1 & 1.9 & 1.0 & 1.2 & 0.8 & 1.7 \\
\midrule
\textit{1919--29} & 2.8 & 4.5 & 5.1 & 3.3 & 3.5 & 3.1 & 4.8 & 9.2 \\
\textit{1929--38} & 2.7 & 3.0 & 3.1 & 2.1 & 1.1 & 0.8 & 0.1 & 5.0 \\
\midrule
\textit{1938--50} & 2.6 & 0.9 & 1.4 & 2.1 & 0.6 & $-$0.1 & 0.8 & 2.4 \\
\textit{1950--68} & 3.9 & 4.9 & 3.6 & 3.1 & 3.6 & 5.2 & 3.9 & 3.4 \\
\textit{1968--89} & 0.9 & 1.8 & 3.6 & 2.5 & 2.4 & 3.0 & 1.7 & 1.5 \\
\textit{1938--89} & 2.4 & 2.7 & 3.1 & 2.6 & 2.4 & 3.0 & 2.3 & 2.4 \\
\midrule
\textit{1989--95} & $-$5.8 & $-$6.6 & $-$0.8 & 1.6 & 1.3 & 2.1 & $-$1.6 & $-$8.0 \\
\textit{1995--2022} & 4.2 & 5.0 & 1.7 & 2.1 & 1.5 & 1.4 & 4.4 & 3.9 \\
\bottomrule
\end{tabular}
\begin{tablenotes}[para,flushleft]
\footnotesize
\textit{Sources:} Estonia: see online appendix. Lithuania: \cite{klimantas2026soviet}. Other series: \cite{mpd2023}, supplemented as documented in the online appendix. \textit{Notes:} compound annual rates between end-point years, \intlnineteen{}. Nordics: Denmark, Finland, Iceland, Norway and Sweden; Western periphery: Finland, Greece, Iceland, Ireland, Italy, Portugal and Spain; Western core: Austria, Belgium, Denmark, France, Germany (West Germany 1950--97), the Netherlands, Norway, Sweden, Switzerland and the UK; ECE (East-Central Europe): Bulgaria, Czechoslovakia, Hungary, Poland, Romania and Yugoslavia; USSR: Soviet Union within 1990 borders (former Soviet Union after 1991).
\end{tablenotes}
\end{threeparttable}
\end{table}


% tab:sectoral
\begin{table}[htbp]
\centering\scriptsize\setlength{\tabcolsep}{3pt}
\caption[Average annual growth of sectoral value-added volume indices, Estonia (per cent)]{\textsc{Average annual growth of sectoral value-added volume indices, Estonia (per cent)}}
\label{tab:sectoral}
\begin{threeparttable}
\begin{tabular}{lrrrrrrcc}
\toprule
Sector & 1919--29 & 1929--38 & 1938--50 & 1950--68 & 1968--89 & 1989--91 & 1938 level & Peak \\
& & & & & & & regained in & year \\
\midrule
Crops & 2.6 & 2.8 & 1.1 & 2.2 & $-$1.1 & $-$8.7 & 1947 & 1968 \\
Livestock & 7.7 & 1.5 & $-$3.9 & 3.7 & 2.0 & $-$8.6 & 1960 & 1987 \\
Forestry and fishing & 4.7 & 0.6 & 0.9 & 7.6 & 1.4 & $-$12.4 & 1947 & 1976 \\
\quad Primary sector & 5.2 & 1.9 & $-$1.3 & 3.1 & 0.6 & $-$8.9 & 1959 & 1986 \\
\midrule
Food & --- & --- & 2.2 & 8.8 & 2.1 & $-$9.7 & 1948 & 1976 \\
Textiles & --- & --- & 4.6 & 9.0 & 0.9 & $-$5.5 & 1949 & 1987 \\
Other light industries & --- & --- & 2.8 & 7.1 & 1.8 & $-$12.4 & 1948 & 1988 \\
Heavy industries & --- & --- & 5.1 & 12.0 & 1.2 & $-$8.6 & 1948 & 1980 \\
\quad Industry & 11.5 & 5.7 & 5.0 & 9.5 & 1.4 & $-$9.1 & 1948 & 1988 \\
\midrule
Construction & 11.7 & 7.6 & 5.4 & 9.3 & 0.9 & $-$9.9 & 1947 & 1977 \\
Transport and communications & 10.7 & 2.0 & 8.0 & 8.9 & 4.9 & $-$4.2 & 1945 & 1989 \\
Trade, accommodation and catering & --- & --- & 1.2 & 9.6 & 4.7 & $-$22.0 & 1949 & 1990 \\
Housing & 1.8 & 6.8 & 1.1 & 2.8 & 2.0 & $-$4.6 & 1949 & 1986 \\
Public services & $-$5.3 & 1.3 & 4.0 & 1.2 & 1.8 & $-$2.5 & 1945 & 1984 \\
\midrule
\textbf{GDP} & 3.3 & 2.9 & 1.9 & 5.4 & 1.7 & $-$8.7 & 1947 & 1989 \\
\bottomrule
\end{tabular}
\begin{tablenotes}[para,flushleft]
\footnotesize
\textit{Sources:} see online appendix. \textit{Notes:} compound annual growth between end-point years of constant-price value-added indices (1938 = 100). Industrial branch indices are available from 1938 only. `1938 level regained': first post-war year (1945 or later) in which the index is at or above its 1938 value; `Peak year': year of the highest index value up to 1990. Public services comprise administration, defence, health and education. Growth in 1938--50 is affected by the wartime trough of 1944, for which Estonian sectoral movements are partly extrapolated from Lithuanian ratios (Section~\ref{sec:limitations}). The 1989--91 column ends in 1991, the splicing year beyond which the final series follows adjusted Maddison Project data.
\end{tablenotes}
\end{threeparttable}
\end{table}


On the eve of the First World War Estonian income per head stood at \ee{1913}~\intlnineteen{}. The war and the fighting of 1918--20 cut it, on the two benchmarks available, which rest on different methods, by more than a fifth, to \ee{1919}~\intlnineteen{} in 1919, a level that was not regained until 1929. Recovery in the 1920s was steady rather than rapid, \gr{1919}{1929} per cent a year, and was carried by the sectors that had suffered most in the war: industry, construction and transport each grew by more than 10 per cent a year from their very low 1919 base, while agriculture, which had been less damaged, grew by \sgr{Primarysectortotal}{1919}{1929} per cent a year (Table~\ref{tab:sectoral}). The Depression was shallow. Income per head fell from \ee{1930}~\intlnineteen{} in 1930 to \ee{1932}~\intlnineteen{} in 1932, a decline of 8 per cent, and had recovered by 1934. The second half of the decade was the fastest phase of interwar growth: with construction growing by \sgr{Construction}{1929}{1938} per cent a year over 1929--38, industry by \sgr{Industrytotal}{1929}{1938} per cent and housing by \sgr{Housing}{1929}{1938} per cent, income per head reached \ee{1938}~\intlnineteen{} in 1938 and an interwar peak of \ee{1940}~\intlnineteen{} in 1940, the year of the Soviet occupation. Over the two decades of independence as a whole, income per head grew by \gr{1919}{1938} per cent a year between 1919 and 1938, but by only \gr{1913}{1938} per cent a year measured from 1913: the recovery from the war absorbed most of the gain.

The Second World War reduced income per head to \ee{1944}~\intlnineteen{} in 1944, a third below the level of 1940, though as Section~\ref{sec:method} explained the 1944 figure is partly extrapolated and should not be pressed. Reconstruction was rapid: GDP regained its 1938 level in \sregain{TotalGVAGDP}, industry in \sregain{Industrytotal} and income per head reached \ee{1950}~\intlnineteen{} in 1950, more than a third above its 1938 level. The result is striking beside Lithuania, where income per head grew by less than 1 per cent a year over 1938--50, and Finland, where it grew by \grx{Finland}{1938}{1950} per cent (Table~\ref{tab:growth}), and two things account for most of it. The first is the denominator: the resident population was 8 per cent smaller in 1950 than in 1938, so that total GDP rose by a quarter where income per head rose by a third. The second is the weight of investment in oil shale, electricity and textiles in the fourth five-year plan, which took industrial output to 1.8 times its 1938 level and transport to 2.5 times by 1950 (Table~\ref{tab:appsec1}). Because the Soviet-period series is anchored in 1990, the 1950 level does not depend on the wartime observations but on growth over 1950--90. Agriculture did not share in this recovery. Crop output was back at its 1938 level by \sregain{Crops}, but livestock output, which had been halved by the war and was then disrupted by collectivisation, did not regain it until \sregain{Livestock}, and the agricultural sector as a whole not until \sregain{Primarysectortotal}, twenty years after the pre-war peak.

The eighteen years from 1950 to 1968 are the most dynamic in the series before 1995. Income per head grew by \gr{1950}{1968} per cent a year, from \ee{1950} to \ee{1968}~\intlnineteen{}, with total GDP growing by \tgr{1950}{1968} per cent a year. Heavy industry grew by \sgr{Heavyindustries}{1950}{1968} per cent a year, textiles by \sgr{Textiles}{1950}{1968} per cent and food processing by \sgr{Food}{1950}{1968} per cent; construction, transport and trade each grew by around 9 per cent a year; and agriculture, now recovering, by \sgr{Primarysectortotal}{1950}{1968} per cent. Growth then slowed sharply and, after 1973, almost ceased. Income per head rose by \gr{1968}{1989} per cent a year between 1968 and 1989, and by only \gr{1973}{1989} per cent a year after 1973; between 1980 and 1989 it rose by \gr{1980}{1989} per cent a year. The sectoral indices show that the slowdown was general. Food processing peaked in \speak{Food}, construction in \speak{Construction}, heavy industry in \speak{Heavyindustries} and the housing stock in \speak{Housing}; only textiles, other light industry, livestock, transport and trade continued to grow into the late 1980s, and industrial output as a whole peaked in \speak{Industrytotal}. The series reaches its Soviet-era maximum of \ee{1989}~\intlnineteen{} in 1989.

The collapse that followed was deeper than anything in the interwar record. Income per head fell by \eeTROUGHpct{} per cent between 1989 and \eeTROUGHyear{}, to \eeTROUGHval~\intlnineteen{}, with the sharpest fall, of 19 per cent, in 1992, the year in which the kroon replaced the rouble and trade with the former Union contracted. Growth resumed in 1995, averaged \gr{1995}{2022} per cent a year to 2022, and carried income per head past its 1989 level in \eeREGAINnineteen{} and to \ee{2022}~\intlnineteen{} (\eeb{2022}~\intltwenty{}) in 2022, second among the Baltic states behind Lithuania.

Table~\ref{tab:previous} compares the new series with the three earlier estimates. For the 1930s the new series and Valge's differ by only a few per cent, which, as explained above, reflects offsetting errors in his domestic-currency total and his conversion ratio rather than agreement in method. For the 1920s the differences are larger and change sign, and the 29 per cent rise that Valge's deflation produces for 1924 becomes a fall of 1 per cent in the new series, in which crops, livestock, industry and transport move by between $-$2 and $+$5 per cent in that year and the small net fall reflects the volatile cement-based construction index (Table~\ref{tab:appsec1}). The series of Klesment, Puur and Valge lies between a quarter and a third below the new one throughout 1960--89, and more than half below in 1950, because of its low UNECE benchmark and the omission of trade and the non-material sphere; its implied growth of income per head over 1950--78, about 4.4 per cent a year, is well above the \gr{1950}{1978} per cent found here, because their index omits the non-material sphere, whose employment-based components grew far more slowly than material production, and weights the four material branches by their shares of net material product, which overstate the share of industry in GDP. The MPD figures, by contrast, lie above the new series, by 7 per cent in 1973 and by 32 per cent in 1989, because they carry modern benchmarks backwards with NMP growth rates that overstate real growth in the final Soviet decades. The new series thus lowers Estonian income in the late Soviet period relative to the MPD, raises it relative to the only previous long series, and replaces both for the years they cover.
